% =====================================================================
%  frontmatter.tex — Bagian Awal Skripsi (format FEB UNDIP)
% =====================================================================
%
%  >>>>>>>>>>>>>>>>>>>>  EDIT DATA DIRI DI SINI  <<<<<<<<<<<<<<<<<<<<<<<
%
\newcommand{\JudulSkripsi}{PENGARUH PRODUK DOMESTIK REGIONAL BRUTO (PDRB)
  DAN INFLASI TERHADAP TINGKAT KEMISKINAN DI PROVINSI JAWA TENGAH
  TAHUN 2019--2023}
\newcommand{\NamaMahasiswa}{Firman Hadi}
\newcommand{\NIM}{12020118130000}
\newcommand{\ProgramStudi}{Ekonomi Pembangunan}
\newcommand{\Jurusan}{Ilmu Ekonomi dan Studi Pembangunan}
\newcommand{\Fakultas}{Ekonomika dan Bisnis}
\newcommand{\Universitas}{Universitas Diponegoro}
\newcommand{\Kota}{Semarang}
\newcommand{\Tahun}{2026}
\newcommand{\DosenPembimbing}{Nama Dosen Pembimbing, S.E., M.Si.}
\newcommand{\NIPPembimbing}{1970xxxxxxxxxx1001}
%
%  >>>>>>>>>>>>>>>>>>>>>>>>>>>>>>>>>>>>>>>>>>>>>>>>>>>>>>>>>>>>>>>>>>>>>>
% =====================================================================

\pagenumbering{roman}
\setcounter{page}{1}

% ------------------------- HALAMAN SAMPUL ----------------------------
\begin{titlepage}
\centering
\thispagestyle{empty}
{\large\bfseries \JudulSkripsi \par}
\vspace{1.5cm}
{\bfseries SKRIPSI\par}
\vspace{0.5cm}
{\normalsize Diajukan sebagai salah satu syarat\\
untuk menyelesaikan Program Sarjana (S1)\\
pada Program Sarjana Fakultas \Fakultas\\ \Universitas \par}
\vspace{1.5cm}
% Ganti berkas gambar/logo-undip.png dengan logo resmi universitas Anda.
\IfFileExists{gambar/logo-undip.png}{%
  \includegraphics[width=4cm]{gambar/logo-undip.png}}{%
  \framebox[4cm]{\rule{0pt}{4cm}\small[ Logo Universitas ]}}
\vspace{1.5cm}

{\normalsize Disusun oleh:\par}
\vspace{0.4cm}
{\bfseries \NamaMahasiswa\par}
{NIM. \NIM\par}
\vspace{1.5cm}
{\bfseries\MakeUppercase{Fakultas \Fakultas}\\
\MakeUppercase{\Universitas}\\
\MakeUppercase{\Kota}\\
\Tahun \par}
\end{titlepage}

% ------------------- HALAMAN PERSETUJUAN SKRIPSI ---------------------
\clearpage
\begin{center}
\textbf{PERSETUJUAN SKRIPSI}
\end{center}
\vspace{1cm}
\begin{singlespace}
\noindent
\begin{tabular}{@{}p{4.5cm}p{0.4cm}p{8.5cm}@{}}
Nama Penyusun        & : & \NamaMahasiswa \\
Nomor Induk Mahasiswa & : & \NIM \\
Fakultas/Jurusan     & : & \Fakultas / \Jurusan \\
Program Studi        & : & \ProgramStudi \\
Judul Skripsi        & : & \JudulSkripsi \\
Dosen Pembimbing     & : & \DosenPembimbing \\
\end{tabular}
\end{singlespace}
\vspace{2cm}
\noindent\hfill \Kota, \ldots\ldots\ldots\ \Tahun

\vspace{0.3cm}
\noindent\hfill Dosen Pembimbing,

\vspace{2.5cm}
\noindent\hfill \underline{\DosenPembimbing}

\noindent\hfill NIP. \NIPPembimbing

% --------------- HALAMAN PENGESAHAN KELULUSAN UJIAN ------------------
\clearpage
\begin{center}
\textbf{PENGESAHAN KELULUSAN UJIAN}
\end{center}
\vspace{0.8cm}
\begin{singlespace}
\noindent
\begin{tabular}{@{}p{4.5cm}p{0.4cm}p{8.5cm}@{}}
Nama Penyusun        & : & \NamaMahasiswa \\
Nomor Induk Mahasiswa & : & \NIM \\
Fakultas/Jurusan     & : & \Fakultas / \Jurusan \\
Judul Skripsi        & : & \JudulSkripsi \\
\end{tabular}
\end{singlespace}
\vspace{0.8cm}

\noindent Telah dinyatakan lulus ujian pada tanggal \ldots\ldots\ldots\ \Tahun.

\vspace{0.6cm}
\noindent Tim Penguji:

\vspace{0.4cm}
\begin{singlespace}
\noindent
\begin{tabular}{@{}p{0.6cm}p{8cm}p{6cm}@{}}
1. & \ldots\ldots\ldots\ldots\ldots\ldots\ldots & (\ldots\ldots\ldots\ldots\ldots) \\[1.2em]
2. & \ldots\ldots\ldots\ldots\ldots\ldots\ldots & (\ldots\ldots\ldots\ldots\ldots) \\[1.2em]
3. & \ldots\ldots\ldots\ldots\ldots\ldots\ldots & (\ldots\ldots\ldots\ldots\ldots) \\
\end{tabular}
\end{singlespace}

% --------------- PERNYATAAN ORISINALITAS SKRIPSI --------------------
\clearpage
\begin{center}
\textbf{PERNYATAAN ORISINALITAS SKRIPSI}
\end{center}
\vspace{0.8cm}
\begin{singlespace}
\noindent Yang bertanda tangan di bawah ini saya, \NamaMahasiswa, menyatakan
bahwa skripsi dengan judul ``\JudulSkripsi'' adalah hasil tulisan saya sendiri.
Dengan ini saya menyatakan dengan sesungguhnya bahwa dalam skripsi ini tidak
terdapat keseluruhan atau sebagian tulisan orang lain yang saya ambil dengan
cara menyalin atau meniru dalam bentuk rangkaian kalimat atau simbol yang
menunjukkan gagasan, pendapat, atau pemikiran dari penulis lain, yang saya akui
seolah-olah sebagai tulisan saya sendiri, dan/atau tidak terdapat bagian atau
keseluruhan tulisan yang saya salin, tiru, atau yang saya ambil dari tulisan
orang lain tanpa memberikan pengakuan penulis aslinya.

\vspace{0.5cm}
\noindent Apabila saya melakukan tindakan yang bertentangan dengan hal tersebut
di atas, baik disengaja maupun tidak, saya menyatakan menarik skripsi yang saya
ajukan sebagai hasil tulisan saya sendiri. Bila kemudian terbukti bahwa saya
melakukan tindakan menyalin atau meniru tulisan orang lain seolah-olah hasil
pemikiran saya sendiri, berarti gelar dan ijazah yang telah diberikan oleh
universitas batal saya terima.
\end{singlespace}
\vspace{1.5cm}
\noindent\hfill \Kota, \ldots\ldots\ldots\ \Tahun

\vspace{0.3cm}
\noindent\hfill Yang membuat pernyataan,

\vspace{2.5cm}
\noindent\hfill \underline{\NamaMahasiswa}

\noindent\hfill NIM. \NIM

% ----------------------------- ABSTRAK ------------------------------
\clearpage
\begin{center}
\textbf{ABSTRAK}
\end{center}
\vspace{0.5cm}
\begin{singlespace}
\noindent Penelitian ini bertujuan menganalisis pengaruh Produk Domestik Regional
Bruto (PDRB) per kapita dan inflasi terhadap tingkat kemiskinan di Provinsi Jawa
Tengah. Data yang digunakan adalah data panel 35 kabupaten/kota selama periode
2019--2023 yang dihimpun dari dua sumber (data utama dan data pendukung) dan
digabungkan melalui kode wilayah, lalu dianalisis dengan metode regresi linear
berganda. Hasil penelitian menunjukkan bahwa PDRB per kapita berpengaruh negatif dan
signifikan terhadap tingkat kemiskinan, sedangkan inflasi berpengaruh positif namun
tidak signifikan. Variabel kontrol Indeks Pembangunan Manusia (IPM) dan akses sanitasi
berpengaruh negatif dan signifikan, sementara tingkat pengangguran berpengaruh positif.
Temuan ini mengindikasikan bahwa peningkatan output ekonomi, kualitas pembangunan
manusia, dan pemenuhan kebutuhan dasar merupakan faktor penting dalam upaya
penanggulangan kemiskinan di Jawa Tengah.

\vspace{0.5cm}
\noindent\textbf{Kata kunci:} kemiskinan, PDRB, inflasi, akses sanitasi, data panel, Jawa Tengah
\end{singlespace}

\clearpage
\begin{center}
\textbf{ABSTRACT}
\end{center}
\vspace{0.5cm}
\begin{singlespace}
\noindent This study aims to analyze the effect of Gross Regional Domestic Product
(GRDP) per capita and inflation on the poverty rate in Central Java Province. The
data used is panel data of 35 regencies/cities during the 2019--2023 period, compiled
from two sources (main and supplementary data) merged by region code, and analyzed
using multiple linear regression. The results show that GRDP per capita has a negative
and significant effect on the poverty rate, while inflation has a positive but
insignificant effect. The control variables, Human Development Index (HDI) and
sanitation access, have a negative and significant effect, while the unemployment rate
has a positive effect. These findings indicate that increasing economic output, the
quality of human development, and the fulfillment of basic needs are important factors
in poverty alleviation efforts in Central Java.

\vspace{0.5cm}
\noindent\textbf{Keywords:} poverty, GRDP, inflation, sanitation access, panel data, Central Java
\end{singlespace}

% -------------------------- KATA PENGANTAR --------------------------
\clearpage
\begin{center}
\textbf{KATA PENGANTAR}
\end{center}
\vspace{0.5cm}
Puji syukur penulis panjatkan ke hadirat Tuhan Yang Maha Esa atas limpahan rahmat
dan karunia-Nya sehingga penulis dapat menyelesaikan skripsi yang berjudul
``\JudulSkripsi''. Skripsi ini disusun sebagai salah satu syarat untuk menyelesaikan
Program Sarjana (S1) pada Fakultas \Fakultas\ \Universitas.

Penulis menyadari bahwa penyusunan skripsi ini tidak lepas dari bantuan berbagai
pihak. Oleh karena itu, penulis menyampaikan terima kasih kepada \ldots
\emph{(tuliskan ucapan terima kasih Anda di sini)}.

Penulis menyadari bahwa skripsi ini masih jauh dari sempurna. Oleh karena itu, kritik
dan saran yang membangun sangat penulis harapkan. Semoga skripsi ini bermanfaat.

\vspace{1cm}
\noindent\hfill \Kota, \ldots\ldots\ldots\ \Tahun

\vspace{0.3cm}
\noindent\hfill Penulis,

\vspace{2cm}
\noindent\hfill \underline{\NamaMahasiswa}

\noindent\hfill NIM. \NIM

% ------------------ DAFTAR ISI / TABEL / GAMBAR ---------------------
\clearpage
{\setstretch{1.2}\tableofcontents}
\clearpage
{\setstretch{1.2}\listoftables}
\clearpage
{\setstretch{1.2}\listoffigures}

% ---------------- Mulai isi utama, nomor halaman arab ---------------
\clearpage
\pagenumbering{arabic}
\setcounter{page}{1}

% Quarto mewajibkan berkas index.qmd sebagai halaman pertama "book", dan
% menerbitkannya sebagai sebuah bab kosong (\chapter{}). Definisi ulang di bawah
% ini "menelan" satu pemanggilan \chapter pertama (milik index.qmd) tanpa
% menghasilkan halaman, lalu mengembalikan perilaku normal — sehingga BAB I
% dimulai tepat dari Pendahuluan.
\let\QuartoFirstChapter\chapter
\renewcommand{\chapter}{\let\chapter\QuartoFirstChapter\setcounter{chapter}{0}}
